\chapter{Отношения и пропорции}
\begin{definition}\term{Отношение} (\emph{ratio}) \(a\) к \(b\) — частное \(\frac{a}{b}\) при \(b \ne  0\); обозначается \(a : b\).\end{definition}
\begin{definition}\term{Пропорция} (\emph{proportion}) — равенство двух отношений: \(\frac{a}{b} = \frac{c}{d}\), где \(b,\allowbreak d \ne  0\).\end{definition}
\begin{definition}\term{Прямая пропорциональность} (\emph{direct proportion}) — зависимость \(y = kx\) с постоянным ненулевым коэффициентом \(k\).\end{definition}
\begin{definition}\term{Обратная пропорциональность} (\emph{inverse proportion}) — зависимость \(y = \frac{k}{x}\) при \(x \ne  0\) с постоянным ненулевым \(k\).\end{definition}
\begin{definition}\term{Масштаб} (\emph{scale}) \(1 : N\) — отношение длины на карте к длине на местности в одинаковых единицах.\end{definition}
\section*{Основное свойство пропорции}
При \(b,\allowbreak d \ne  0\) равенство \(\frac{a}{b} = \frac{c}{d}\) равносильно \(a \cdot  d = b \cdot  c\).\par
\section*{Правила применения}
При прямой пропорциональности увеличение \(x\) в несколько раз увеличивает \(y\) во столько же раз.\par
При обратной пропорциональности положительных величин увеличение \(x\) в несколько раз уменьшает \(y\) во столько же раз.\par
Деление \(S\) в отношении \(m : n\) при \(m,\allowbreak n > 0\): части равны \(S \cdot  \frac{m}{m + n}\) и \(S \cdot  \frac{n}{m + n}\).\par
{\small\textit{Пример: \(60\) разделить в отношении \(2 : 3\) — это \(24\) и \(36\), поскольку одна часть равна \(\frac{60}{5} = 12\).}}\par
\Needspace{7\baselineskip}\section*{Задачи}
\begin{enumerate}[label={},leftmargin=3.2em,itemsep=10pt]
\Needspace{3\baselineskip}\item[\textbf{5.1.}] Сократи отношение \(18 : 24\).
\Needspace{3\baselineskip}\item[\textbf{5.2.}] Реши пропорцию \(\frac{x}{15} = \frac{4}{5}\).
\Needspace{3\baselineskip}\item[\textbf{5.3.}] Реши пропорцию \(\frac{7}{x} = \frac{21}{18}\).
\Needspace{3\baselineskip}\item[\textbf{5.4.}] Раздели \(64\) в отношении \(3 : 5\).
\Needspace{3\baselineskip}\item[\textbf{5.5.}] Для \(4\) порций нужно \(300\) г крупы. Сколько нужно для \(10\) порций?
\Needspace{3\baselineskip}\item[\textbf{5.6.}] \(5\) одинаковых тетрадей стоят \(180\) рублей. Сколько стоят \(8\)?
\Needspace{3\baselineskip}\item[\textbf{5.7.}] На карте масштаба \(1 : 100 000\) расстояние \(3,5\) см. Найди расстояние на местности в километрах.
\Needspace{3\baselineskip}\item[\textbf{5.8.}] \(6\) рабочих выполняют работу за \(10\) дней. За сколько дней справятся \(15\) рабочих при той же производительности?
\Needspace{3\baselineskip}\item[\textbf{5.9.}] Поездка при скорости \(60\) км/ч занимает \(3\) часа. Сколько займёт при \(90\) км/ч на том же пути?
\Needspace{3\baselineskip}\item[\textbf{5.10.}] Смешивают сок и воду в отношении \(2 : 3\). Сколько каждого компонента нужно для \(1,5\) л смеси?
\Needspace{3\baselineskip}\item[\textbf{5.11.}] При \(y = 3 x\) найди \(y\) для \(x = 4\). Как изменится \(y\) при удвоении \(x\)?
\Needspace{3\baselineskip}\item[\textbf{5.12.}] При \(y = \frac{24}{x}\) найди \(y\) для \(x = 6\). Как изменится \(y\) при удвоении положительного \(x\)?
\Needspace{3\baselineskip}\item[\textbf{5.13.}] Из жизни. Для \(6\) порций супа нужно \(450\) г картофеля. Сколько нужно для \(10\) таких же порций?
\Needspace{3\baselineskip}\item[\textbf{5.14.}] Из жизни. \(4\) одинаковых насоса заполняют бассейн за \(9\) часов. За сколько часов его заполнят \(6\) таких насосов, работающих одновременно?
\Needspace{3\baselineskip}\item[\textbf{5.15.}] Из жизни. Для напитка сироп смешивают с водой в отношении \(1 : 7\). Сколько сиропа и воды нужно для \(2\) л напитка?
\Needspace{3\baselineskip}\item[\textbf{5.16.}\starmark] Рабочие одинаковой производительности: \(6\) человек выполняют работу за \(8\) дней по \(5\) часов. За сколько дней справятся \(4\) человека по \(6\) часов?
\Needspace{3\baselineskip}\item[\textbf{5.17.}\starmark] В смеси отношение сиропа к воде \(2 : 5\). После добавления \(3\) л воды стало \(1 : 4\). Сколько литров смеси было сначала?
\Needspace{3\baselineskip}\item[\textbf{5.18.}\starmark] Автомобиль ехал одинаковое время со скоростями \(60\) и \(90\) км/ч. Найди среднюю скорость. А если одинаковыми были расстояния?
\Needspace{3\baselineskip}\item[\textbf{5.19.}] Пусть \(b,d>0\) и \(\frac ab=\frac cd=k\). Докажи, что \(\frac{a+c}{b+d}=k\).
\Needspace{3\baselineskip}\item[\textbf{5.20.}\starmark] Пусть \(a,b,c,d>0\) и \(\frac ab<\frac cd\). Докажи, что \(\frac ab<\frac{a+c}{b+d}<\frac cd\).
\end{enumerate}
